\documentclass[../../main.tex]{subfiles}
\begin{document}

\begin{center}
  {\fontsize{48pt}{58pt}\selectfont \textbf{2004}\normalsize\textbf{年度}}
\end{center}

\begin{table}[htb]
  \centering
  \caption{出題テーマと難易度}
  \begin{tabular}{|c|p{9cm}|c|}
    \hline
    \textbf{問題番号} & \textbf{テーマ} & \textbf{難易度} \\
    \hline
    第1問 & 放物線上の正三角形と座標幾何 & \\
    \hline
    第2問 & 平方数の下位桁と整除 & \\
    \hline
    第3問 & 円内を転がる円板が描く曲線と面積 & \\
    \hline
    第4問 & 三次関数の反復と方程式の実数解の個数 & \\
    \hline
    第5問 & 二球の和集合の体積とグラフ & \\
    \hline
    第6問 & 板の裏返し操作と確率 & \\
    \hline
  \end{tabular}
\end{table}

\begin{table}[htb]
  \centering
  \caption{解答}
  \renewcommand{\arraystretch}{1.4}
  \begin{tabular}{|c|c|p{9cm}|}
    \hline
    \textbf{問題番号} & \textbf{小問} & \textbf{答え} \\
    \hline
    第1問 & -- & $a=\dfrac{18}{5}$ \\
    \hline
    第2問 & (1) & $b=0$または$4$（証明問題） \\
    \cline{2-3}
    & (2) & 平方数は$10000$で割り切れる（証明問題） \\
    \hline
    第3問 & -- & 二つの面積は$\dfrac{108\pi}{5}$，$\dfrac{392\pi}{5}$ \\
    \hline
    第4問 & (1) & $|a|>2$で1個，$|a|=2$で2個，$|a|<2$で3個 \\
    \cline{2-3}
    & (2) & $|a|>2$で1個，$|a|=2$で5個，$|a|<2$で9個 \\
    \cline{2-3}
    & (3) & $f_n(x)=0$の実数解は$3^n$個（証明問題） \\
    \hline
    第5問 & (1) & $0<r\le2$で
    $V=\pi\left(-\dfrac{r^3}{12}+r+\dfrac{4}{3}\right)$，
    $r\ge2$で$V=\dfrac{8\pi}{3}$（グラフは解答参照） \\
    \cline{2-3}
    & (2) & $r=1.5$（小数第1位まで） \\
    \hline
    第6問 & (1) & $\dfrac{7}{27}$ \\
    \cline{2-3}
    & (2) & $\dfrac{1}{2}\left(\dfrac{1}{3}\right)^n
    +\dfrac{1}{4}\left(-\dfrac{1}{3}\right)^n+\dfrac{1}{4}$ \\
    \hline
  \end{tabular}
\end{table}

\paragraph{試験形態}
6問．試験時間は要確認．

\end{document}
